\chapter{Regolarizzazione: dropout, batch normalization e oltre}
\label{cap:regolarizzazione}
\fonti{\rif{2026-09-30}{30/09}, \rif{2026-10-06}{06/10}, \rif{2026-10-20}{20/10}}{2-foundament-deep-learning, pp.\ 50--51;
2a-dropout\_batchnorm, pp.\ 1--7; notebook FFNN}

\section{Overfitting e regolarizzazione}
Una rete va in \textbf{overfitting} quando è (quasi) perfetta solo sui dati di training: la loss di training continua a scendere,
anche verso 0, mentre quella di \textbf{validation}, dopo essere scesa, \textbf{risale} (learning curve e diagnosi nel
cap.~\ref{cap:addestramento}). Il modello ha memorizzato il dataset invece di imparare.

\begin{definizione}[regolarizzazione]
Insieme di tecniche che aggiungono \textbf{vincoli o penalità} al processo di apprendimento per ridurre l'overfitting e migliorare la
\textbf{generalizzazione}. È cruciale nel deep learning perché le reti neurali hanno una grande capacità di \textbf{memorizzare} i
dati di training. Le penalità si aggiungono alla loss:
\[
\mathcal{L} = \mathcal{L}_{std} + \text{penalty term}.
\]
\end{definizione}

Le slide raggruppano le tecniche secondo la ``leva'' su cui agiscono:
\begin{center}
\small
\begin{tabularx}{\textwidth}{l l X}
\toprule
\textbf{Si agisce su\ldots} & \textbf{Tecnica} & \textbf{Idea}\\
\midrule
Pesi & L1 (Lasso), L2 (Ridge) / weight decay & penalità nella loss sul valore dei pesi\\
Architettura & Dropout & spegnere a caso neuroni durante il training\\
Processo di training & Early stopping & fermarsi quando la validation peggiora\\
Dati & Batch normalization & normalizzare gli input dei layer con le statistiche del mini-batch\\
Dati & Noise injection & rumore casuale su input, pesi o attivazioni\\
\bottomrule
\end{tabularx}
\end{center}

\section{Regolarizzazione sui pesi: L1 e L2}
\begin{itemize}
  \item \textbf{L1 (Lasso)}: penalità proporzionale al \textbf{valore assoluto} dei pesi, $\lambda\sum_i|\phi_i|$. Spinge alcuni pesi
        \textbf{a zero} $\Rightarrow$ sparsità, una sorta di selezione delle feature: parti diverse della rete si specializzano.
  \item \textbf{L2 (Ridge)}: penalità proporzionale al \textbf{quadrato} dei pesi, $\lambda\sum_i\phi_i^2$. Scoraggia i
        \textbf{pesi grandi} e distribuisce l'importanza su più feature.
  \item $\lambda$ è la \textbf{forza della regolarizzazione} (un iperparametro).
  \item In PyTorch la L2 si ottiene con il parametro \code{weight\_decay} dell'ottimizzatore (es.\ \code{AdamW}, vedi
        cap.~\ref{cap:addestramento}); un valore tipico di partenza indicato nelle linee guida di progettazione è $0{,}01$
        (cap.~\ref{cap:ffnn}).
\end{itemize}
\begin{integrazione}[sul ``rumore'']
A lezione L1/L2 sono state descritte come ``aggiungere un po' di rumore'' all'aggiornamento dei pesi. È un'intuizione:
tecnicamente sono \textbf{penalità} deterministiche nella loss che impediscono ai pesi di adattarsi perfettamente al training
set. L'aggiunta di rumore vera e propria è la \en{noise injection} (sezione~\ref{sec:noise-injection}).
\end{integrazione}

\section{Dropout}
\begin{definizione}[dropout]
Durante ogni iterazione di training si ``spegne'' (si pone a zero l'output di) un sottoinsieme \textbf{casuale} di neuroni.
I neuroni non possono diventare troppo dipendenti gli uni dagli altri $\Rightarrow$ la rete impara \textbf{feature più robuste}.
In predizione si usano \textbf{tutti} i neuroni.
\end{definizione}

\subsection{Cosa succede in training}
\begin{itemize}
  \item A ogni forward pass di training, l'output di ogni neurone è posto a 0 con probabilità $p$.
  \item La scelta è \textbf{indipendente per ogni campione} del mini-batch, e nuova a ogni mini-batch: ogni campione è elaborato da una
        \textbf{sotto-rete diversa}. Un neurone spento ora può essere acceso dopo (non si ``congela'' una volta per tutte).
  \item Gli output sopravvissuti sono \textbf{scalati di} $1/(1-p)$, così il valore atteso dell'output resta uguale con e senza dropout.
  \item Obiettivo: i neuroni non possono fare affidamento su altri neuroni specifici $\Rightarrow$ feature più robuste, meno overfitting.
  \item Il dropout \textbf{non rimuove} neuroni né pesi: azzera temporaneamente degli output.
\end{itemize}

\subsection{Dropout e backpropagation}
\begin{itemize}
  \item Forward, backward e aggiornamento avvengono \textbf{una volta per mini-batch}; il gradiente usato è la \textbf{media} sul mini-batch.
  \item Tutti i pesi vengono aggiornati, ma \textbf{ogni neurone impara solo dai campioni in cui era acceso}: si riduce la capacità di
        apprendimento e ogni neurone vede campioni diversi.
\end{itemize}

\subsection{Training vs validation/test}
\begin{center}
\small
\begin{tabularx}{\textwidth}{X X}
\toprule
\textbf{Training} (\code{model.train()}) & \textbf{Validation/Test} (\code{model.eval()})\\
\midrule
neuroni casuali spenti, nuova maschera a ogni passaggio & tutti i neuroni usati, nessuna maschera\\
sopravvissuti scalati di $1/(1-p)$ & nessuno scaling: il dropout diventa l'identità\\
output casuale: stesso input, output diversi & output deterministico, usa i pesi appresi\\
output 2, $p=0{,}5$ $\to$ 0 oppure 4 (in media 2) & output 2 $\to$ 2 (stesso valore atteso)\\
\bottomrule
\end{tabularx}
\end{center}
\begin{attenzione}[\code{no\_grad} non basta]
\code{torch.no\_grad()} disattiva solo i gradienti: è \code{model.eval()} che spegne il dropout. Se si valuta senza
\code{model.eval()}, le predizioni restano casuali.
\end{attenzione}

\subsection{Quanto dropout?}
Con $p=0{,}5$ metà dei neuroni del layer ha output 0 durante un passo di training. Secondo il docente $0{,}5$ è probabilmente troppo:
conviene provare valori tra $0{,}2$ e $0{,}5$ (stessa indicazione nelle tabelle di progettazione di FFNN, RNN e GNN:
cap.~\ref{cap:ffnn}, \ref{cap:rnn}, \ref{cap:cnn-gnn}).

\begin{integrazione}[validation loss più bassa della training con dropout]
Nel notebook d'esempio, con dropout $0{,}5$ la validation loss risulta \textbf{più bassa} di quella di training. È normale: in
training il dropout spegne neuroni e peggiora la loss, in validation la rete lavora al completo.
\end{integrazione}

\section{Batch normalization}
Un layer tra due layer che \textbf{ri-normalizza} i dati mentre attraversano la rete: normalizza gli input di uno o più layer per
\textbf{stabilizzare e velocizzare} il training. Riduce l'overfitting rendendo il training \textbf{meno sensibile all'inizializzazione}
dei pesi e si può applicare anche \textbf{prima} della funzione di attivazione.

\subsection{Cosa succede in training}
Per ogni \textbf{feature}, sull'output del layer precedente (\code{N} campioni $\times$ \code{F} feature):
\begin{enumerate}
  \item media $\mu_B$ e varianza $\sigma^2_B$ \textbf{del mini-batch corrente} (statistiche calcolate lungo il batch);
  \item standardizzazione: $\hat x = \dfrac{x-\mu_B}{\sqrt{\sigma^2_B+\epsilon}}$ $\Rightarrow$ media 0, varianza 1;
  \item scala e traslazione: $y = \gamma\hat x + \beta$, con $\gamma, \beta$ \textbf{appresi} (parametri della rete, aggiornati dalla
        backpropagation): la rete impara quanto conta la normalizzazione.
\end{enumerate}
\begin{itemize}
  \item Batch diversi $\Rightarrow$ $\mu, \sigma$ diversi: l'output di un campione \textbf{dipende dagli altri campioni del batch}
        (da qui il nome).
  \item A ogni mini-batch si aggiornano anche \textbf{media e varianza mobili} (\en{running}), \textbf{non apprese} ma calcolate:
        \[ \text{running} = 0{,}9\cdot\text{running} + 0{,}1\cdot\text{statistica del batch} \]
\end{itemize}

\subsection{Training vs validation/test}
\begin{itemize}
  \item \textbf{Training}: usa $\mu, \sigma$ del mini-batch corrente e aggiorna le running statistics.
  \item \textbf{Validation/Test}: usa le \textbf{running mean/variance memorizzate} (congelate). La BN \textbf{non sparisce}: diventa una
        scala e traslazione fissa. Senza, il layer successivo riceverebbe in validation dati con una distribuzione diversa da quella
        vista in training.
\end{itemize}
\begin{esempio}[stesso valore, output diversi]
Valore $x=6$ di una feature, $\gamma=1$, $\beta=0$:
\begin{center}
\small
\begin{tabular}{llcc}
\toprule
\textbf{Fase} & \textbf{Batch} & \textbf{Statistiche} & \textbf{Output}\\
\midrule
Training & A $=[2,4,6,8]$ & $\mu=5$, $\sigma\approx2{,}24$ & $+0{,}45$\\
Training & B $=[6,10,12,14]$ & $\mu=10{,}5$, $\sigma\approx2{,}96$ & $-1{,}52$\\
Validation/Test & qualsiasi & running $\mu=5$, $\sigma\approx2{,}24$ & $+0{,}45$\\
\bottomrule
\end{tabular}
\end{center}
Stesso input, output diversi in training: questo ``rumore di batch'' dà un lieve effetto di \textbf{regolarizzazione}.
\end{esempio}

\begin{nota}[Dropout e batch norm in PyTorch]
\begin{center}
\small
\begin{tabularx}{\linewidth}{l X X l}
\toprule
& \code{model.train()} & \code{model.eval()} & \textbf{Appreso in backprop}\\
\midrule
Dropout & maschera casuale per campione, scala $1/(1-p)$ & identità: tutti i neuroni, nessuno scaling & niente ($p$ fisso)\\
BatchNorm & $\mu,\sigma$ del batch + aggiorna le running stats & running $\mu,\sigma$ (congelate) & $\gamma, \beta$\\
\bottomrule
\end{tabularx}
\end{center}
\end{nota}

\section{Early stopping}
Si monitora la loss (o la performance) sul \textbf{validation set}: quando inizia a peggiorare mentre quella di training continua a
migliorare, si \textbf{interrompe} il training, prima che il modello si adatti troppo ai dati di training. Si usa il modello
\textbf{prima} dell'overfitting (all'epoch $t$, non $t+10$).
\begin{itemize}
  \item \code{min\_delta}: miglioramento minimo della validation loss per considerarlo significativo.
  \item \code{patience}: numero massimo di epoch consecutive senza miglioramento significativo prima di fermarsi (5, 10, 20\ldots).
  \item Più lunga la patience, più a lungo si addestra; più grande il min\_delta, più forte deve essere il miglioramento.
        Una richiesta troppo severa (calo di 0,5 in 2 epoch) ferma tutto subito.
  \item \code{best\_model\_state}: si \textbf{salvano i pesi del modello migliore} e li si ricarica alla fine. Senza, si
        restituirebbe il modello dell'ultima epoch (già in overfitting), oppure bisognerebbe ri-addestrare da zero fino all'epoch
        migliore: ore perse con dataset grandi.
\end{itemize}
\begin{lstlisting}
import copy
best_val, wait = float('inf'), 0
patience, min_delta = 20, 1e-4
best_state = None
for epoch in range(num_epochs):
    ...                                   # training + validation, calcola val_loss
    if val_loss < best_val - min_delta:   # miglioramento significativo
        best_val, wait = val_loss, 0
        best_state = copy.deepcopy(model.state_dict())
    else:
        wait += 1
        if wait >= patience:
            break                         # stop
model.load_state_dict(best_state)         # ripristina il modello migliore
\end{lstlisting}

\section{Noise injection}\label{sec:noise-injection}
Si aggiunge \textbf{rumore casuale} a input, pesi o attivazioni durante il training per rendere il modello più robusto.
È citata nelle slide ma non la useremo nel corso.

\section{Regolarizzazione nel codice}
In PyTorch dropout e batch norm sono \textbf{layer} come gli altri, definiti in \code{\_\_init\_\_}; la L2 è un parametro
dell'ottimizzatore.
\begin{lstlisting}
class RegularizedNN(nn.Module):
    def __init__(self, input_size, num_classes, p=0.5):
        super().__init__()
        self.fc1 = nn.Linear(input_size, 128)
        self.fc2 = nn.Linear(128, 128)
        self.bn2 = nn.BatchNorm1d(128)        # batch normalization
        self.drop = nn.Dropout(p)             # dropout: si definisce come un layer
        self.fc3 = nn.Linear(128, num_classes)
        self.relu = nn.ReLU()
    def forward(self, x):
        x = self.relu(self.fc1(x))
        x = self.relu(self.bn2(self.fc2(x)))  # BN prima della prossima trasformazione
        x = self.drop(x)
        return self.fc3(x)

optimizer = optim.AdamW(model.parameters(), lr=1e-3, weight_decay=1e-2)  # L2
\end{lstlisting}

\begin{esempio}[cosa mostra il notebook d'esempio]
Sul dataset giocattolo della lezione del \rif{2026-10-06}{06/10} (pipeline completa nel cap.~\ref{cap:ffnn}):
\begin{itemize}
  \item \textbf{early stopping} (patience 20, min\_delta piccolo): il training si ferma molto prima;
  \item \textbf{dropout 0,5}: curva più stabile; la validation loss è più bassa di quella di training;
  \item \textbf{batch norm + weight decay}: curva molto più rumorosa. Le prestazioni restano buone solo perché l'esempio è facile:
        ``se vedo una validation loss così nel vostro report, penso che non abbiate valutato la soluzione''.
\end{itemize}
\end{esempio}

\section{Quando la regolarizzazione non basta}
\begin{attenzione}
Se l'overfitting è forte perché la rete è enorme e memorizza i dati, i rimedi ``leggeri'' non bastano: riduci layer e
neuroni, oppure ripensa preprocessing e architettura.
\end{attenzione}
Tra gli iperparametri, il \code{weight\_decay} è quello da toccare \textbf{se c'è overfitting}; il learning rate resta il primo su cui
lavorare (cap.~\ref{cap:ffnn}).

\begin{sintesi}
\begin{itemize}
  \item La regolarizzazione aggiunge vincoli o penalità all'apprendimento ($\mathcal{L} = \mathcal{L}_{std} + \text{penalty}$) per
        combattere l'overfitting.
  \item \textbf{Pesi}: L1 (sparsità, pesi a zero) e L2 (pesi piccoli; in PyTorch \code{weight\_decay}).
  \item \textbf{Dropout}: in training azzera ogni output con probabilità $p$ (maschera per campione, scala $1/(1-p)$); in
        \code{eval()} è l'identità. Valori tipici $0{,}2$--$0{,}5$.
  \item \textbf{Batch norm}: standardizza ogni feature con $\mu_B,\sigma_B$ del mini-batch, poi $\gamma\hat x+\beta$ appresi; in
        \code{eval()} usa le running statistics congelate. Stabilizza il training, lieve effetto regolarizzante.
  \item \textbf{Early stopping}: \code{patience} e \code{min\_delta} sulla validation loss, e si ripristinano i pesi migliori.
  \item Ricorda \code{model.train()}/\code{model.eval()}: è ciò che accende e spegne dropout e statistiche di batch.
\end{itemize}
\end{sintesi}
